% !PRINTED 395
\weberpage{395}
\begin{center}
\S.\,127.

The Pell Equation.
\end{center}

If $\omega$ is a reduced number belonging to the discriminant $D$, then
by the continued fraction development one obtains
\weq{1}{\omega = \dfrac{P_n \omega_n + P_{n-1}}{Q_n \omega_n + Q_{n-1}},}
and $\omega_n = \omega$ holds always when $n$ is a multiple of the
number of terms $\nu$ of the period. But then it follows from (1):
\weq{2}{Q_n \omega^2 = (P_n - Q_{n-1}) \omega + P_{n-1}.}
Now let $\omega$ satisfy the quadratic equation \S.\,122, (5):
\weq{3}{c \omega^2 = a + b \omega, \quad D = b^2 + 4 a c,}
and from this, since $a$, $b$, $c$ are without common divisor,
equation (2) must follow by multiplication with an integer.
If we denote this integer by $u$, we thus have
\weq{4}{Q_n = u c, \quad P_{n-1} = u a, \quad P_n - Q_{n-1} = u b.}
If we further set
\weq{5}{P_n + Q_{n-1} = t,}
so that $t$ too is a positive integer, then from (4) and (5) it follows that
\[
P_n = \frac{t + u b}{2}, \quad Q_n = u c,
\]
\weq{6}{P_{n-1} = u a, \quad Q_{n-1} = \frac{t - u b}{2},}
and hence, having regard to the equation
\[
P_n Q_{n-1} - Q_n P_{n-1} = (-1)^n\text{:}
\]
\weq{7}{t^2 - D u^2 = (-1)^n 4.}
Equation (7) is called the \gesperrt{Pell equation}. The
problem is this: for a given $D$, to find all its integral
solutions $t$, $u$. Since, if $t$, $u$ satisfy the equation,
$\pm t$, $\pm u$ also satisfy it, we may restrict ourselves to the determination
of its positive solutions.

The formulae (4), (5) give us an infinite number of such
solutions; we have only to choose for $n$ an arbitrary multiple of the
period number $\nu$, and for $u$ the greatest common
divisor of $Q_n$, $P_{n-1}$, $P_n - Q_{n-1}$, or also simply the number
%CONT>>
